% ==============================================================================
% 第 34 章：小结 (Chapter 34: Summary)
% ==============================================================================

\chapter{第六部分小结}
\label{sec:ch34_summary}

本书最后一部分系统性地探讨了\term{累加器}（accumulators）的设计之道。累加器是一种在数据结构遍历过程中持续收集、维系并传递上下文知识的优雅而强大的计算机制。为函数配备累加器，不仅能够修复程序中灾难性的时间与空间性能瓶颈，更能够在生成型递归中彻底消除路径循环、保证算法严格终止并产生正确的结果。

从本部分的学习与操练中，你应该领悟并牢固掌握“两个半”核心设计法则：

\begin{recipebox}{累加器设计的核心法则}
\begin{enumerate}
  \item \textbf{第一步：敏锐识别引入累加器的契机}。
  常规的自顶向下遍历在由表及里、从当前片段迈向下一个片段时，天然会“遗忘”先前的部分实参信息。如果你发现这些遗失的上下文知识能够大幅简化后续函数的设计、消除二次遍历或避免回路死循环，就应当果断切换至累加器模板。

  \item \textbf{第二步：严谨表述累加器不变量（核心跃迁）}。
  这是掌握累加器编程最关键、也是最需要智识投入的一步。必须用精确的自然语言或代数等式说明：累加器究竟以何种数据结构收集了何种上下文知识。在绝大多数场景下，\textbf{累加器不变量精准描述了主函数的初始实参与辅助函数当前实参之间的差量（Difference）关系}。

  \item \textbf{第三步：由不变量直接导出实现三要素（水到渠成）}。
  有了精确的不变量，后续的编码实现便有了清晰的数学指引：
  \begin{enumerate}
    \item 准确确定累加器的初始启动值 $a_0$；
    \item 准确确定在深入递归的单步操作中如何更新累加器值；
    \item 准确确定如何在基准条件（递归终点）或递归分支中充分压榨并利用累加器所维系的知识。
  \end{enumerate}
\end{enumerate}
\end{recipebox}

累积上下文知识的思维模式在整个计算机科学与软件工程领域无处不在，并以极其多元的形态贯穿于各类编程范式之中。在以 ISL+、Scheme、Racket、Haskell 为代表的函数式编程语言中，它通过函数参数的显式单向流动来实现，天然具备无副作用、易于形式化推导以及享受尾调用优化（Tail-Call Optimization, TCO）转化为常数空间执行的卓越特性。

而在主流命令式编程语言（如 C、Java、Python、Go）中，程序员往往在 \code{for}、\code{while} 等原生循环构造中通过\textbf{变量赋值}（assignment）来累加状态，因为命令式循环语句本身通常不返回值。然而，在设计这些命令式循环程序时，其思维心法与本书传授的累加器不变量设计配方\textbf{本质上完全相通}——所谓“循环不变量”（Loop Invariant），正是此处累加器不变量在命令式时空中的孪生镜像！严谨推导不变量，永远是编写出无 Bug、高性能高质量软件的终极护城河。
